\section{Spectral Analysis Framework}
\label{sec:method}

\subsection{Spectral Optimization via Orthogonalized Updates}

Muon~\cite{jordan2024muon,liu2025muon} optimizes 2D weight matrices by orthogonalizing the gradient before applying the update, effectively performing a Newton-like step in the spectral domain.
Given a gradient matrix $G \in \mathbb{R}^{m \times n}$ for a weight matrix $W$, the goal is to compute an approximation to $(GG^\top)^{-1/2} G$, which projects the gradient onto the Stiefel manifold of orthonormal row matrices.
Rather than forming this matrix inverse explicitly, Muon uses Newton--Schulz iteration to approximate it efficiently.
Five iterations of a quintic Newton--Schulz polynomial orthogonalize $G$ (starting from $X_0 = G/\|G\|_F$; iteration formula, coefficients, and convergence analysis in Appendix~\ref{sec:muon_background}).
The orthogonalized gradient serves as the update direction, combined with Nesterov momentum ($\beta = 0.95$).

This procedure captures inter-parameter curvature structure that first-order methods miss.
SGD updates along the raw gradient direction; AdamW applies element-wise scaling via running second-moment estimates, which amounts to a diagonal approximation to the curvature.
Muon's orthogonalization operates on the full gradient matrix, providing a non-diagonal second-order approximation related to K-FAC~\cite{martens2015kfac}, but without explicitly maintaining Kronecker-factored estimates.
As described in Section~\ref{sec:experimental_setup}, Muon is applied exclusively to 2D weight matrices (convolution kernels reshaped to 2D and linear layers), while remaining parameters are handled by a separate AdamW instance (Appendix Table~\ref{tab:muon_routing} details the parameter routing).

\subsection{Layer-wise Spectral Analysis Framework}

Global accuracy differences between optimizers do not reveal \emph{where} in the network one optimizer outperforms another.
We introduce two complementary analysis tools that localize optimizer effects to individual layers.

\paragraph{Weight Swap Protocol.}
Given a Muon-trained model with weights $\theta^{\text{Muon}} = \{W_l^M\}$ and an AdamW-trained model with weights $\theta^{\text{AdamW}} = \{W_l^A\}$, we construct a hybrid model $\theta^{\text{Muon} \setminus l \leftarrow \text{AdamW}_l}$ by replacing the weights of a single layer $l$ in the Muon model with the corresponding AdamW weights, keeping all other layers at their Muon-trained values.
We define the per-layer Muon benefit:
\begin{equation}
    \Delta a_l = \mathrm{acc}(\theta^{\text{Muon}}) - \mathrm{acc}(\theta^{\text{Muon} \setminus l \leftarrow \text{AdamW}_l}),
    \label{eq:layer_benefit}
\end{equation}
where $\mathrm{acc}(\cdot)$ denotes validation accuracy.
$\Delta a_l > 0$ indicates that Muon's training at layer $l$ outperforms AdamW; $\Delta a_l < 0$ indicates the reverse.
The beneficial layer ratio, the fraction of layers where $\Delta a_l > 0$, aggregates this across all eligible layers.

\paragraph{Spectral Correlation Analysis.}
To connect the weight swap results with spectral structure, we measure how each optimizer shapes the \emph{stable rank} of each layer's weight matrix:
\begin{equation}
    \mathrm{sr}(W_l) = \frac{\|W_l\|_F^2}{\sigma_1(W_l)^2},
    \label{eq:spectral_response}
\end{equation}
where $\sigma_1(W_l)$ is the largest singular value.
Stable rank captures the effective dimensionality of the weight matrix: higher values indicate a more uniform singular value distribution, meaning the matrix spreads its representational capacity across more directions.
We choose stable rank over alternatives (effective rank, condition number, spectral gap) because it is bounded, well-defined for rank-deficient matrices, and permits direct cross-architecture comparison without architecture-specific normalization.
For convolutional layers, we reshape the 4D kernel tensor to a 2D matrix before computing $\mathrm{sr}(W_l)$.
We compute the per-layer stable rank divergence $\Delta \mathrm{sr}_l = \mathrm{sr}(W_l^{\text{Muon}}) - \mathrm{sr}(W_l^{\text{AdamW}})$: positive values indicate Muon spreads spectral energy more uniformly (higher effective rank), negative values indicate Muon concentrates energy into fewer directions.
We evaluate the Spearman rank correlation between $\Delta \mathrm{sr}_l$ and $\Delta a_l$ (Eq.~\ref{eq:layer_benefit}); results are reported in Section~\ref{sec:spectral_analysis}.

\subsection{Evaluation Metrics and Statistical Protocol}

We report best validation accuracy (early-stopping checkpoint) throughout, in percentage points (pp), to decouple optimizer comparison from convergence speed.
Key experiments use 3--4 seeds (from $\{42, 0, 1, 2\}$; Appendix Table~\ref{tab:multi_seed}), reporting mean $\pm$ std ($\text{ddof}=1$).
Independent 50-epoch runs with cosine annealing ($T_{\max} = 50$) and the same per-pair learning rates test gap persistence without schedule-length confounding.